% Options for packages loaded elsewhere
\PassOptionsToPackage{unicode}{hyperref}
\PassOptionsToPackage{hyphens}{url}
\documentclass[
  11pt,
]{article}
\usepackage{xcolor}
\usepackage[margin=1in]{geometry}
\usepackage{amsmath,amssymb}
\setcounter{secnumdepth}{-\maxdimen} % remove section numbering
\usepackage{iftex}
\ifPDFTeX
  \usepackage[T1]{fontenc}
  \usepackage[utf8]{inputenc}
  \usepackage{textcomp} % provide euro and other symbols
\else % if luatex or xetex
  \usepackage{unicode-math} % this also loads fontspec
  \defaultfontfeatures{Scale=MatchLowercase}
  \defaultfontfeatures[\rmfamily]{Ligatures=TeX,Scale=1}
\fi
\usepackage{lmodern}
\ifPDFTeX\else
  % xetex/luatex font selection
\fi
% Use upquote if available, for straight quotes in verbatim environments
\IfFileExists{upquote.sty}{\usepackage{upquote}}{}
\IfFileExists{microtype.sty}{% use microtype if available
  \usepackage[]{microtype}
  \UseMicrotypeSet[protrusion]{basicmath} % disable protrusion for tt fonts
}{}
\makeatletter
\@ifundefined{KOMAClassName}{% if non-KOMA class
  \IfFileExists{parskip.sty}{%
    \usepackage{parskip}
  }{% else
    \setlength{\parindent}{0pt}
    \setlength{\parskip}{6pt plus 2pt minus 1pt}}
}{% if KOMA class
  \KOMAoptions{parskip=half}}
\makeatother
\usepackage{color}
\usepackage{fancyvrb}
\newcommand{\VerbBar}{|}
\newcommand{\VERB}{\Verb[commandchars=\\\{\}]}
\DefineVerbatimEnvironment{Highlighting}{Verbatim}{commandchars=\\\{\}}
% Add ',fontsize=\small' for more characters per line
\newenvironment{Shaded}{}{}
\newcommand{\AlertTok}[1]{\textcolor[rgb]{1.00,0.00,0.00}{\textbf{#1}}}
\newcommand{\AnnotationTok}[1]{\textcolor[rgb]{0.38,0.63,0.69}{\textbf{\textit{#1}}}}
\newcommand{\AttributeTok}[1]{\textcolor[rgb]{0.49,0.56,0.16}{#1}}
\newcommand{\BaseNTok}[1]{\textcolor[rgb]{0.25,0.63,0.44}{#1}}
\newcommand{\BuiltInTok}[1]{\textcolor[rgb]{0.00,0.50,0.00}{#1}}
\newcommand{\CharTok}[1]{\textcolor[rgb]{0.25,0.44,0.63}{#1}}
\newcommand{\CommentTok}[1]{\textcolor[rgb]{0.38,0.63,0.69}{\textit{#1}}}
\newcommand{\CommentVarTok}[1]{\textcolor[rgb]{0.38,0.63,0.69}{\textbf{\textit{#1}}}}
\newcommand{\ConstantTok}[1]{\textcolor[rgb]{0.53,0.00,0.00}{#1}}
\newcommand{\ControlFlowTok}[1]{\textcolor[rgb]{0.00,0.44,0.13}{\textbf{#1}}}
\newcommand{\DataTypeTok}[1]{\textcolor[rgb]{0.56,0.13,0.00}{#1}}
\newcommand{\DecValTok}[1]{\textcolor[rgb]{0.25,0.63,0.44}{#1}}
\newcommand{\DocumentationTok}[1]{\textcolor[rgb]{0.73,0.13,0.13}{\textit{#1}}}
\newcommand{\ErrorTok}[1]{\textcolor[rgb]{1.00,0.00,0.00}{\textbf{#1}}}
\newcommand{\ExtensionTok}[1]{#1}
\newcommand{\FloatTok}[1]{\textcolor[rgb]{0.25,0.63,0.44}{#1}}
\newcommand{\FunctionTok}[1]{\textcolor[rgb]{0.02,0.16,0.49}{#1}}
\newcommand{\ImportTok}[1]{\textcolor[rgb]{0.00,0.50,0.00}{\textbf{#1}}}
\newcommand{\InformationTok}[1]{\textcolor[rgb]{0.38,0.63,0.69}{\textbf{\textit{#1}}}}
\newcommand{\KeywordTok}[1]{\textcolor[rgb]{0.00,0.44,0.13}{\textbf{#1}}}
\newcommand{\NormalTok}[1]{#1}
\newcommand{\OperatorTok}[1]{\textcolor[rgb]{0.40,0.40,0.40}{#1}}
\newcommand{\OtherTok}[1]{\textcolor[rgb]{0.00,0.44,0.13}{#1}}
\newcommand{\PreprocessorTok}[1]{\textcolor[rgb]{0.74,0.48,0.00}{#1}}
\newcommand{\RegionMarkerTok}[1]{#1}
\newcommand{\SpecialCharTok}[1]{\textcolor[rgb]{0.25,0.44,0.63}{#1}}
\newcommand{\SpecialStringTok}[1]{\textcolor[rgb]{0.73,0.40,0.53}{#1}}
\newcommand{\StringTok}[1]{\textcolor[rgb]{0.25,0.44,0.63}{#1}}
\newcommand{\VariableTok}[1]{\textcolor[rgb]{0.10,0.09,0.49}{#1}}
\newcommand{\VerbatimStringTok}[1]{\textcolor[rgb]{0.25,0.44,0.63}{#1}}
\newcommand{\WarningTok}[1]{\textcolor[rgb]{0.38,0.63,0.69}{\textbf{\textit{#1}}}}
\setlength{\emergencystretch}{3em} % prevent overfull lines
\providecommand{\tightlist}{%
  \setlength{\itemsep}{0pt}\setlength{\parskip}{0pt}}
\usepackage{bookmark}
\IfFileExists{xurl.sty}{\usepackage{xurl}}{} % add URL line breaks if available
\urlstyle{same}
% fallback for those not using the hyperref driver hyperxmp:
\makeatletter
\@ifundefined{xmpquote}{\newcommand{\xmpquote}[1]{#1}}{}
\makeatother
\hypersetup{
  hidelinks,
  pdfcreator={LaTeX via pandoc}}

\author{}
\date{}

\usepackage{fvextra}
\DefineVerbatimEnvironment{Highlighting}{Verbatim}{commandchars=\\\{\},breaklines,fontsize=\small}
\begin{document}

\section{Mathematical contract and exact acceptance
derivation}\label{mathematical-contract-and-exact-acceptance-derivation}

\subsection{Scope and provenance}\label{scope-and-provenance}

Hypra is a \textbf{partial native Python/SymPy implementation} for exact
regulator expansions. This chapter gives detailed proofs for three
illustrative families and the P/Q/J4 acceptance calculation, followed by
the general integer and half-integer kernels. These examples are not the
package's full supported scope; consult the \href{scope.md}{capability
map} for current coverage and limits. The native backend is the default
and \textbf{requires no Mathematica installation or setup}. Unsupported
cases must fail explicitly rather than be presented as complete
expansions.

HypExp {[}1{]} expands generalized hypergeometric functions around
integer-valued parameters. HypExp 2 {[}2{]} adds certain half-integer
classes. Hypra does not yet match their full coverage. An optional
external HypExp adapter is retained as an experimental,
\textbf{unverified} alternative; it is neither required by the native
implementation nor evidence that a real Wolfram/HypExp run succeeded.

The calculations below independently derive the requested example.
Earlier private project material was context only: no private code or
integral lookup table is reproduced or imported here. The reproducible
evidence is Hypra's public Python validation source, not a private file
reference.

\subsubsection{Public implementation
excerpt}\label{public-implementation-excerpt}

\href{https://github.com/lalmei/hypra/blob/4b2904aa1163f9e47f2f497cd0b9a9be5c719b59/scripts/validate_integral.py\#L20-L45}{\texttt{derive\_moments()},
lines 20--45} derives the five logarithmic moments using convergent
sums, beta-function coefficients and a dilogarithm reflection identity.
The excerpt below is copied verbatim apart from indentation from
\href{https://github.com/lalmei/hypra/blob/4b2904aa1163f9e47f2f497cd0b9a9be5c719b59/scripts/validate_integral.py\#L26-L38}{lines
26--38} of the pinned published revision \texttt{4b2904a}. Here
\texttt{s} is SymPy, and \texttt{a} and \texttt{t} are formal expansion
variables.

\begin{Shaded}
\begin{Highlighting}[]
\NormalTok{n }\OperatorTok{=}\NormalTok{ s.Symbol(}\StringTok{"n"}\NormalTok{, integer}\OperatorTok{=}\VariableTok{True}\NormalTok{, positive}\OperatorTok{=}\VariableTok{True}\NormalTok{)}
\NormalTok{a, t }\OperatorTok{=}\NormalTok{ s.symbols(}\StringTok{"a t"}\NormalTok{)}
\NormalTok{log\_gamma\_ratio }\OperatorTok{=}\NormalTok{ s.expand(}
    \BuiltInTok{sum}\NormalTok{((}\OperatorTok{{-}}\DecValTok{1}\NormalTok{) }\OperatorTok{**}\NormalTok{ k }\OperatorTok{*}\NormalTok{ s.zeta(k) }\OperatorTok{*}\NormalTok{ (a}\OperatorTok{**}\NormalTok{k }\OperatorTok{+}\NormalTok{ t}\OperatorTok{**}\NormalTok{k }\OperatorTok{{-}}\NormalTok{ (a }\OperatorTok{+}\NormalTok{ t) }\OperatorTok{**}\NormalTok{ k) }\OperatorTok{/}\NormalTok{ k }\ControlFlowTok{for}\NormalTok{ k }\KeywordTok{in} \BuiltInTok{range}\NormalTok{(}\DecValTok{2}\NormalTok{, }\DecValTok{5}\NormalTok{))}
\NormalTok{)}
\CommentTok{\# Beta(a, 1+t) = exp(log\_gamma\_ratio)/a. Taking d\_a d\_t\^{}2}
\CommentTok{\# means twice the a\^{}2*t\^{}2 coefficient before division by a.}
\NormalTok{beta\_mixed }\OperatorTok{=}\NormalTok{ s.expand(log\_gamma\_ratio }\OperatorTok{+}\NormalTok{ log\_gamma\_ratio}\OperatorTok{**}\DecValTok{2} \OperatorTok{/} \DecValTok{2}\NormalTok{)}
\NormalTok{b }\OperatorTok{=}\NormalTok{ s.simplify(}\DecValTok{2} \OperatorTok{*}\NormalTok{ beta\_mixed.coeff(a, }\DecValTok{2}\NormalTok{).coeff(t, }\DecValTok{2}\NormalTok{))}
\NormalTok{zeta2\_from\_series }\OperatorTok{=}\NormalTok{ s.summation(n}\OperatorTok{**{-}}\DecValTok{2}\NormalTok{, (n, }\DecValTok{1}\NormalTok{, s.oo))}
\CommentTok{\# Li2(1{-}w) = zeta(2){-}Li2(w){-}log(w)log(1{-}w), and}
\CommentTok{\# integral Li2(w)log(1{-}w) dw/w = {-}Li2(1)\^{}2/2.}
\NormalTok{c }\OperatorTok{=}\NormalTok{ s.simplify(}\OperatorTok{{-}}\NormalTok{(zeta2\_from\_series}\OperatorTok{**}\DecValTok{2}\NormalTok{) }\OperatorTok{/} \DecValTok{2} \OperatorTok{{-}}\NormalTok{ b)}
\end{Highlighting}
\end{Shaded}

The mixed coefficient \texttt{b} gives
\(\int_0^1 \log(w)\log^2(1-w)\,dw/w\); \texttt{c} is the moment
containing \(\log(1-w)\operatorname{Li}_2(1-w)\). The mathematical
justification for extracting these coefficients is developed below. The
subsequent
\href{https://github.com/lalmei/hypra/blob/4b2904aa1163f9e47f2f497cd0b9a9be5c719b59/scripts/validate_integral.py\#L39-L45}{returned
moment dictionary} uses exact convergent zeta sums, and
\href{https://github.com/lalmei/hypra/blob/4b2904aa1163f9e47f2f497cd0b9a9be5c719b59/scripts/validate_integral.py\#L48-L89}{\texttt{derive\_j4()}}
combines the resulting moments into the quadratic and cubic regulator
coefficients. No fitted or hard-coded J4 coefficient is inserted. This
Python validation does not claim that an external Mathematica/HypExp
installation was run.

\subsection{Conventions}\label{conventions}

\begin{itemize}
\tightlist
\item
  The expansion variable is a dimensionless regulator \(\epsilon\),
  expanded at zero.
\item
  \texttt{through=k} means every Laurent term through and
  \textbf{including} \(\epsilon^k\), plus a remainder
  \(O(\epsilon^{k+1})\). Pole terms must not be discarded.
\item
  Inputs and coefficients are exact. A binary floating-point
  approximation is not a substitute for a rational parameter.
\item
  The derivation in this document uses \textbf{real \(0<w<1\)},
  principal real logarithms and classical polylogarithms. Continuation
  outside this interval requires a branch prescription; the formulas
  must not be silently reused across cuts.
\item
  An HPL object retains its complete index word and argument. Under the
  usual real alphabet convention, \(f_0(t)=1/t\), \(f_1(t)=1/(1-t)\) and
  \(f_{-1}(t)=1/(1+t)\), with \(H_0(z)=\log z\), \(H_1(z)=-\log(1-z)\),
  and \(H_{-1}(z)=\log(1+z)\). Compact nonzero index \(m\) denotes
  \(|m|-1\) zeros followed by \(\operatorname{sgn}(m)\), not a different
  kernel. Complex continuation, extended alphabets, and endpoint
  regularization are separate matters; a wrapper must not invent their
  conventions.
\end{itemize}

\subsection{Native all-order HPL recursion for an integer
family}\label{native-all-order-hpl-recursion-for-an-integer-family}

The native implementation is based on a differential equation, not
stored acceptance coefficients. Let

\[F(z,\epsilon)={}_2F_1(1,b\epsilon;1+c\epsilon;z)
=1+\sum_{k\geq1}\epsilon^k f_k(z),\]

where \(b,c\) are regulator-independent exact constants. Its
argument-series coefficient \(a_n=(b\epsilon)_n/(1+c\epsilon)_n\)
satisfies \((n+c\epsilon)a_n=(n-1+b\epsilon)a_{n-1}\) for \(n\geq1\).
Summing this identity gives

\[z(1-z)\partial_zF+\epsilon(c-bz)F=c\epsilon,\qquad F(0,\epsilon)=1.\]

At order one \(f_1'=b/(1-z)\), hence \(f_1=bH_1(z)\). At every higher
order,

\[f_k'(z)=\left[-\frac{c}{z}+\frac{b-c}{1-z}\right]f_{k-1}(z),
\qquad f_k(0)=0.\]

With \(\mathcal I_a\) denoting the operation of prepending letter \(a\)
to an HPL word, this is the constructive recurrence

\[\boxed{f_1=bH_1,\qquad
f_k=-c\,\mathcal I_0 f_{k-1}+(b-c)\,\mathcal I_1 f_{k-1}\quad(k\geq2).}\]

Every word ends in 1, so its origin value is zero and these integrals
require no divergent origin regularization. Equivalently, a word of
length \(k\) ending in 1, with \(n_0\) zero letters and \(n_1\) one
letters, has coefficient \(b(-c)^{n_0}(b-c)^{n_1-1}\). This proves every
order by induction and uniqueness of the regular solution, not only the
first two acceptance orders. It covers \(P\) with \((b,c)=(-1,-2)\) and
\(Q\) with \((b,c)=(-2,-2)\) and argument \(1-w\).

The derivation starts on \(|z|<1\) near \(\epsilon=0\). For real
\(0<z<1\) it uses the unambiguous real HPL branch. Analytic continuation
follows only with matching HPL path/branch conventions; no branch cut
crossing is implicit. The offline recurrence tests compare orders 1
through 5, at argument powers 1 through 6, against exact Pochhammer
quotients for four choices of \((b,c)\). The proof, not that finite test
sample, establishes the all-order recurrence.

\subsection{Two narrowly supported half-integer
families}\label{two-narrowly-supported-half-integer-families}

These identities do \textbf{not} establish general half-integer HypExp 2
coverage. Let \(t=\sqrt z\) with real \(0<z<1\).

\subsubsection{Even-binomial family}\label{even-binomial-family}

For \(a\) independent of the regulator,

\[G(z,\epsilon)={}_2F_1(a\epsilon,\tfrac12+a\epsilon;\tfrac12;z)
=\frac{(1-t)^{-2a\epsilon}+(1+t)^{-2a\epsilon}}2.\]

To prove it, take the even terms of the binomial series and use
\((2A)_{2n}=2^{2n}(A)_n(A+1/2)_n\) and \((2n)!=2^{2n}n!(1/2)_n\), with
\(A=a\epsilon\). Thus

\[[\epsilon^k]G=\frac{(-2a)^k}{2k!}
\left[\log^k(1-t)+\log^k(1+t)\right],\quad k\geq0.\]

The expression is even in \(t\) and has a regular \(z=0\) limit. Its
derivation from the binomial identity is exact, not a fitted expansion.
No cancellation of upper and lower hypergeometric parameters at nonzero
regulator was used.

\subsubsection{Half-integer telescoping
family}\label{half-integer-telescoping-family}

For \(A=1/2+b\epsilon\),

\[T(z,\epsilon)={}_2F_1(1,A;1+A;z)
=\sum_{n\geq0}\frac{A}{n+A}z^n.\]

Define

\[S_k(z)=\sum_{n\geq0}\frac{z^n}{(n+1/2)^k}
=\frac{2^{k-1}}t\left[\operatorname{Li}_k(t)-\operatorname{Li}_k(-t)\right].\]

The equality follows by extracting odd terms of the polylogarithm
series. Since \(A/(n+A)=1-n/(n+A)\), direct geometric expansion of its
denominator gives

\[T_0=\frac{\log(1+t)-\log(1-t)}{2t},\qquad
[\epsilon^k]T=(-1)^{k-1}b^k
\left[S_k-\frac12S_{k+1}\right]\quad(k\geq1).\]

This includes the \(n=0\) term correctly: it contributes 1 at order zero
and zero at every higher order. For each fixed finite \(b\), the
expansion is local in \(\epsilon\), before the nearest denominator pole.
At \(z=0\) the displayed quotients have removable limits (\(T_0=1\),
higher coefficients zero), whereas \(z=1\) is divergent. An
implementation must handle those endpoints explicitly, not evaluate a
spurious \(0/0\) or silently return a finite value at 1.

\subsection{P and Q from their defining
series}\label{p-and-q-from-their-defining-series}

Set

\[P(w,\epsilon)={}_2F_1(1,-\epsilon;1-2\epsilon;w),\qquad
Q(w,\epsilon)={}_2F_1(1,-2\epsilon;1-2\epsilon;1-w).\]

The upper parameter \(-2\epsilon\) in \(Q\) is \textbf{not} the lower
parameter \(1-2\epsilon\); no upper/lower cancellation applies.

For \(n\geq1\), using \((1)_n/n!=1\),

\[\frac{(-\epsilon)_n}{(1-2\epsilon)_n}
=-\frac{\epsilon}{n}\left[1+\epsilon
\left(H_{n-1}+\frac2n\right)+O(\epsilon^2)\right].\]

This follows by writing the rising factorials as products and
differentiating \(\sum_{j=1}^{n-1}\log(1-\epsilon/j)
-\sum_{j=1}^{n}\log(1-2\epsilon/j)\). On \(|w|<1\),

\[\sum_{n\geq1}\frac{w^n}{n}=-\log(1-w),\quad
\sum_{n\geq1}\frac{w^n}{n^2}=\operatorname{Li}_2(w),\quad
\sum_{n\geq1}\frac{H_{n-1}w^n}{n}=\frac12\log^2(1-w).\]

The last identity follows by differentiating both sides, using
\(\sum_{n\geq1}H_{n-1}w^{n-1}=-\log(1-w)/(1-w)\), and the value zero at
\(w=0\). Consequently

\[P=1+\epsilon\log(1-w)
+\epsilon^2\left[-2\operatorname{Li}_2(w)-\frac12\log^2(1-w)\right]
+O(\epsilon^3).\]

For \(Q\) there is an exact telescoping coefficient:

\[\frac{(-2\epsilon)_n}{(1-2\epsilon)_n}=-\frac{2\epsilon}{n-2\epsilon}.
\qquad
Q=1-\sum_{j\geq1}(2\epsilon)^j\operatorname{Li}_j(1-w).\]

The double-series interchange is justified for fixed \(0<w<1\) and
\(|2\epsilon|<1\) by absolute convergence. Hence

\[Q=1+2\epsilon\log w-4\epsilon^2\operatorname{Li}_2(1-w)
+O(\epsilon^3).\]

The finite-\(n\) coefficient identities are checked exactly by the
offline tests; they are not numerically fitted.

\subsection{Why the J4 expansion may be
integrated}\label{why-the-j4-expansion-may-be-integrated}

Define

\[J_4(\epsilon)=\int_0^1
w^{-1-2\epsilon}(1-w)^{-\epsilon}(P-1)(Q-1)\,dw.\]

Fix a sufficiently small closed complex disk \(|\epsilon|\leq r<1/4\).
All hypergeometric denominator parameters avoid their singular values
there. The exact \(Q\) coefficients imply
\(|Q-1|\leq C_r|\epsilon||\log w|\) near \(w=0\) and
\(|Q-1|\leq C_r|\epsilon|(1-w)\) near \(w=1\). For \(P\), its
coefficient can be written

\[-\epsilon\frac{\Gamma(1-2\epsilon)}{\Gamma(1-\epsilon)}
\frac{\Gamma(n-\epsilon)}{\Gamma(n+1-2\epsilon)}.\]

Uniform gamma-ratio asymptotics on this compact regulator disk give
\(|p_n|\leq C_r|\epsilon|n^{r-1}\). Therefore \(P-1=O_r(|\epsilon|w)\)
near zero and \(P-1=O_r(|\epsilon|(1-w)^{-r})\) near one. The full
integrand is bounded by constants times

\[w^{-2r}|\log w|\quad(w\to0),\qquad
(1-w)^{1-2r}\quad(w\to1),\]

both integrable. Dominated holomorphic parameter integration and
Cauchy's formula on a smaller disk permit differentiation under this
integral to any fixed finite order. Thus the Taylor coefficients
computed below are legitimate integrals; the integrated remainder is
\(O(\epsilon^4)\). This proof is specific to \(J_4\), \textbf{not}
permission to expand each unsubtracted singular term in a larger NNLO
integral separately.

\subsection{Deriving all needed moments, not importing a
table}\label{deriving-all-needed-moments-not-importing-a-table}

All integrals in this subsection have measure \(dw/w\) on \((0,1)\).
Write \(L=\log w\), \(M=\log(1-w)\) and \(\zeta_k=\zeta(k)\). Expanding
the logarithm or dilogarithm and integrating their absolutely integrable
terms gives

\[\int LM=\sum_{n\geq1}\frac1{n^3}=\zeta_3,\qquad
\int L\operatorname{Li}_2(w)=-\sum_{n\geq1}\frac1{n^4}=-\zeta_4,\]

\[\int L^2M=-2\sum_{n\geq1}\frac1{n^4}=-2\zeta_4.\]

Here \(\int_0^1w^{n-1}\log w\,dw=-1/n^2\) and
\(\int_0^1w^{n-1}\log^2w\,dw=2/n^3\). The absolute sums converge.

For \(B=\int LM^2\), use \(\mathrm B(a,1+t)-1/a\), analytic at
\((a,t)=(0,0)\), and take \(\partial_a\partial_t^2\). The relevant
generating function is

\[\mathrm B(a,1+t)=\frac1a\exp\left\{
\sum_{k\geq2}\frac{(-1)^k\zeta_k}{k}
[a^k+t^k-(a+t)^k]\right\}.\]

The \(a^2t^2\) coefficient of the exponential is
\(-3\zeta_4/2+\zeta_2^2/2=-\zeta_4/4\). After division by \(a\) and the
two \(t\) derivatives, \(B=-\zeta_4/2\). Differentiation is justified
either after subtracting \(1/a\) or directly for the mixed derivative,
whose logarithmic integrand is integrable.

Finally, with \(C=\int M\operatorname{Li}_2(1-w)\), the real reflection
identity \(\operatorname{Li}_2(w)+\operatorname{Li}_2(1-w)=\zeta_2-LM\)
gives

\[C=-\zeta_2^2+\frac12\zeta_2^2-B
=-\frac12\zeta_2^2-B=-\frac34\zeta_4.\]

We used \(d\operatorname{Li}_2(w)=-M\,dw/w\) and its endpoint values;
all boundary terms are finite.

\subsection{Assemble the exact
coefficients}\label{assemble-the-exact-coefficients}

Expand the weight as \(w^{-1}[1+\epsilon(-2L-M)+O(\epsilon^2)]\).
Multiplying the derived series gives

\[[\epsilon^2]J_4=\int 2LM=2\zeta_3,\]

\[[\epsilon^3]J_4=\int[-4L\operatorname{Li}_2(w)-3LM^2
-4M\operatorname{Li}_2(1-w)-4L^2M]
=\left(4+\frac32+3+8\right)\zeta_4=\frac{33}{2}\zeta_4.\]

Thus the requested result is consistent:

\[\boxed{J_4=2\zeta(3)\epsilon^2+\frac{11\pi^4}{60}\epsilon^3
+O(\epsilon^4).}\]

The script creates the polynomial product and integrates each monomial
using the independently derived moments. It rejects any unproved
monomial rather than supplying a guessed value.

\subsection{Laurent prefactors and order
budgeting}\label{laurent-prefactors-and-order-budgeting}

If a prefactor \(A\) starts at \(\epsilon^{-p}\), an analytic factor
\(F\) is generally needed through order \(k+p\) to obtain \(AF\) through
order \(k\). For example, \(\Gamma(\epsilon)^2F\) needs \(F\) through
\(\epsilon^{k+2}\), not merely \(\epsilon^k\). Multiplication, exact
cancellations and the final truncation must be performed on the
\textbf{whole expression}. If a backend cannot establish the
valuation/order budget or its result still contains an unexpanded
hypergeometric function, a regulator derivative or a regulator-dependent
coefficient, it must fail rather than label the result complete.
Assumptions needed to establish a valuation or cancel a parameter are
part of the mathematical input.

\subsection{Validation and limits}\label{validation-and-limits}

\begin{Shaded}
\begin{Highlighting}[]
\ExtensionTok{uv}\NormalTok{ run python scripts/validate\_integral.py}
\ExtensionTok{uv}\NormalTok{ run python scripts/validate\_integral.py }\AttributeTok{{-}{-}latex}
\ExtensionTok{uv}\NormalTok{ run python scripts/validate\_integral.py }\AttributeTok{{-}{-}numeric}
\ExtensionTok{uv}\NormalTok{ run pytest tests/test\_mathematics.py}
\end{Highlighting}
\end{Shaded}

The optional numerical check evaluates the original integral using
contiguous representations for \(P-1\) and \(Q-1\) to avoid catastrophic
subtraction near the endpoints. It uses negative regulator samples, for
which endpoint limits are zero, and introduces \textbf{no endpoint
cutoff}. Convergence of the residual scaled by \(\epsilon^4\) is
secondary support only, not the analytic proof and not a real-backend
validation. No global continuation or arbitrary-HPL simplifier is
established by these examples.

\subsection{References}\label{references}

\begin{enumerate}
\def\labelenumi{\arabic{enumi}.}
\tightlist
\item
  T. Huber and D. Maître, \emph{HypExp, a Mathematica package for
  expanding hypergeometric functions around integer-valued parameters},
  \href{https://arxiv.org/abs/hep-ph/0507094}{arXiv:hep-ph/0507094}.
\item
  T. Huber and D. Maître, \emph{HypExp 2, Expanding Hypergeometric
  Functions about Half-Integer Parameters},
  \href{https://arxiv.org/abs/0708.2443}{arXiv:0708.2443}.
\end{enumerate}

The abstract pages were checked for these scope statements. No source
code from these packages is redistributed with this derivation.

\subsection{Native general-parameter
kernels}\label{native-general-parameter-kernels}

The implementation now constructs canonical integer coefficients by
strict nested-sum algebra and their HPL generating functions. Matched
half-integer parameters use even-index projections and colored strict
sums. Parameter shifts act with the hypergeometric differential operator

\[
 L=\theta\prod_j(\theta+b_j-1)-z\prod_i(\theta+a_i),
 \qquad \theta=z\frac{d}{dz}.
\]

For example, raising an upper parameter applies
\(F(a_i+1)=(\theta+a_i)F/a_i\), and lowering a lower parameter applies
\(F(b_j-1)=(\theta+b_j-1)F/(b_j-1)\). The inverse shifts follow
polynomial division of the same operator. Each zero-base divisor
increases the canonical regulator depth before any truncation. Apparent
Laurent poles must cancel exactly below the proven pole bound. An
identically zero shift divisor is resolved by a single formal slope
perturbation only when its finite limit can be established
algebraically.

For all-half bases at rank \(p\),

\[
 L_0=(\theta-\tfrac12)^{p-1}B,
 \qquad B=(1-z)\theta-\tfrac z2.
\]

The order of these operators is essential. Set \(F=(1-z)^{-1/2}Y\),
\(t=i\sqrt{z/(1-z)}\), so \(z=t^2/(t^2-1)\). The gauge-conjugated
operators are

\[
 \widetilde\theta=\frac{t(1-t^2)}2\partial_t-\frac{t^2}2,
 \quad \widetilde D=\frac{t(1-t^2)}2\partial_t-\frac{1+t^2}2,
 \quad \widetilde B=\frac t2\partial_t.
\]

Their origin-normalized inverses are

\[
 \widetilde D^{-1}g(t)=\frac{t}{1-t^2}\int_0^t\frac{2g(u)}{u^2}\,du,
 \qquad
 \widetilde B^{-1}g(t)=\int_0^t\frac{2g(u)}u\,du.
\]

At each regulator order apply \(\widetilde D^{-1}\) \(p-1\) times and
then \(\widetilde B^{-1}\). The initial coefficient is \(Y_0=1\);
correction sources and solutions are \(O(t^2)\), making these integrals
convergent. The odd homogeneous solution of \(\widetilde D\) is excluded
by analyticity in \(z\), and the constant homogeneous solution of
\(\widetilde B\) is fixed by \(F(0)=1\). Only HPL letters \(0,\pm1\) are
required. For real \(0<z<1\), continuation follows the positive
imaginary \(t\) axis with principal square roots. This does not assert
arbitrary complex-argument continuation or permission to interchange
integrals and limits.

These are independent Python-native mathematical implementations
inspired by the two cited HypExp papers. The detailed capability map and
resource restrictions are in \texttt{scope.md}; canonical class coverage
is not a claim of complete HypExp/HPL API parity.

\end{document}
